\documentclass[12pt,a4paper]{article}

% ---- Encoding & Font ----
\usepackage{iftex}
\ifPDFTeX
  \usepackage[utf8]{inputenc}
  \usepackage[T1]{fontenc}
\else
  \usepackage{fontspec}  % XeTeX/LuaTeX (tectonic): Unicode nativo (·, —, ¿, ª, §)
\fi
\usepackage[spanish]{babel}
\usepackage{csquotes}

% ---- Page layout ----
\usepackage[top=2.5cm, bottom=2.5cm, left=2.5cm, right=2.5cm]{geometry}
\usepackage{setspace}
\onehalfspacing

% ---- Math ----
\usepackage{amsmath, amssymb, amsthm}
\usepackage{mathtools}
\usepackage{physics}
\usepackage{esint}  % \oiint

% ---- Graphics ----
\usepackage{graphicx}
\usepackage{tikz}
\usepackage{float}
\usepackage{caption}

% ---- Tables ----
\usepackage{array}
\usepackage{booktabs}
\usepackage{tabularx}
\usepackage{colortbl}

% ---- Colours ----
\usepackage{xcolor}
\definecolor{notebg}{HTML}{FFF3CD}
\definecolor{noteborder}{HTML}{FFC107}
\definecolor{boxred}{HTML}{F8D7DA}
\definecolor{boxredborder}{HTML}{F5C6CB}

% ---- Boxes ----
\usepackage{mdframed}
\usepackage{tcolorbox}
\tcbuselibrary{most}

\newtcolorbox{keybox}[1][]{
  colback=blue!5,
  colframe=blue!70!black,
  arc=4pt,
  boxrule=1pt,
  fonttitle=\bfseries,
  title={#1}
}

\newtcolorbox{warnbox}[1][]{
  colback=boxred,
  colframe=boxredborder,
  coltitle=black!75,
  arc=4pt,
  boxrule=1pt,
  fonttitle=\bfseries,
  title={#1}
}

\newtcolorbox{notebox}[1][]{
  colback=notebg,
  colframe=noteborder,
  coltitle=black!75,
  arc=4pt,
  boxrule=1pt,
  fonttitle=\bfseries,
  title={#1}
}

% ---- Hyperlinks & TOC ----
\usepackage[hidelinks]{hyperref}
\usepackage{bookmark}

% ---- Enumerate ----
\usepackage{enumitem}

% ---- Miscellaneous ----
\usepackage{icomma}
\usepackage{siunitx}
\usepackage{textcomp}
\usepackage[bottom]{footmisc}

% ---- Title ----
\title{\textbf{Clase 10: Capacitores Esférico y Cilíndrico, Asociación \\
        Serie/Paralelo, Energía Almacenada y Dieléctricos}\\[0.3em]
        \large{Capacitancia por geometría, capacidad equivalente,
              energía en el campo y polarización}}
\author{Transcripto y expandido --- Curso Física III (OpenFING)}
\date{}

\begin{document}

\maketitle
\thispagestyle{empty}
\newpage

\tableofcontents
\newpage

% ===================================================================
\section{Demostración experimental: equipotenciales}
% ===================================================================

Dos estudiantes muestran una \textbf{cuba electrolítica} (agua como
conductor débil) a 5\,V, midiendo el potencial con un tester sobre curvas
en papel cuadriculado. Se verifica que el potencial es constante sobre
cada equipotencial y que $\mathbf{E}$ es perpendicular a ellas.
Configuraciones:

\begin{itemize}[nosep]
  \item \textbf{Dos líneas de carga} ($\approx$ placas): equipotenciales
        rectas paralelas.
  \item \textbf{Aros concéntricos} ($\approx$ carga central + cascarón):
        equipotenciales circulares.
  \item \textbf{Dipolo}: plano bisector a potencial medio; esferas cerca
        de cada carga.
  \item \textbf{Plano + carga puntual}: rectas cerca del plano,
        curvándose cerca de la carga.
\end{itemize}

\begin{notebox}[Simulación numérica (relajación)]
  Fijar el potencial en los bordes y reemplazar iterativamente cada punto
  por el \textbf{promedio de sus vecinos} hasta converger.
\end{notebox}

% ===================================================================
\section{Capacitor de placas esféricas}
% ===================================================================

Cascarón conductor interno de radio $A$ ($+Q$) y externo de radio $B$
($-Q$). Receta general: campo $\to \Delta V \to C$.

\textbf{Campo (Gauss)}, esfera de radio $r$ ($A<r<B$):
\[
E\cdot 4\pi r^2 = \frac{Q}{\varepsilon_0}
\;\Longrightarrow\;
E = \frac{Q}{4\pi\varepsilon_0 r^2}.
\]

\textbf{Diferencia de potencial} (segmento radial):
\[
|\Delta V| = \int_A^B \frac{Q}{4\pi\varepsilon_0 r^2}\,dr
           = \frac{Q}{4\pi\varepsilon_0}\left(\frac{1}{A}-\frac{1}{B}\right).
\]

\begin{keybox}[Capacitor esférico]
  \[
  \boxed{C = \frac{4\pi\varepsilon_0}{\dfrac{1}{A}-\dfrac{1}{B}}}
  \]
\end{keybox}

% ===================================================================
\section{Capacitor de placas cilíndricas}
% ===================================================================

Cilindros concéntricos de radios $A$ ($+Q$) y $B$ ($-Q$), largo
$L \gg B$ (se desprecian bordes).

\textbf{Campo (Gauss)}, cilindro de radio $r$ y altura $H$; las tapas no
contribuyen y la carga encerrada es $Q\,H/L$:
\[
E\cdot 2\pi r H = \frac{Q\,H/L}{\varepsilon_0}
\;\Longrightarrow\;
E = \frac{Q}{2\pi\varepsilon_0 L\, r} = \frac{\lambda}{2\pi\varepsilon_0 r},
\]
idéntico al de una línea infinita con $\lambda = Q/L$.

\textbf{Diferencia de potencial:}
\[
|\Delta V| = \int_A^B \frac{Q}{2\pi\varepsilon_0 L\, r}\,dr
           = \frac{Q}{2\pi\varepsilon_0 L}\ln\!\frac{B}{A}.
\]

\begin{keybox}[Capacitor cilíndrico]
  \[
  \boxed{C = \frac{2\pi\varepsilon_0 L}{\ln(B/A)}}
  \]
\end{keybox}

\begin{notebox}[Bordes]
  Las aproximaciones valen lejos de los bordes; cilindros cortos tienen
  efectos de borde adicionales.
\end{notebox}

% ===================================================================
\section{Condensadores en serie y en paralelo}
% ===================================================================

Ciertos grupos de condensadores equivalen a uno solo. Se supone
equilibrio electrostático y cables ideales (cada tramo conductor es un
equipotencial).

\subsection{Paralelo}

Comparten $|\Delta V_{AB}|$; la carga se reparte: $Q = Q_1 + Q_2$. Con
$Q_i = |\Delta V_{AB}|\,C_i$:

\begin{keybox}[Paralelo]
  \[
  \boxed{C_{\text{eq}} = C_1 + C_2}
  \qquad (\text{se suman las capacitancias})
  \]
\end{keybox}

\subsection{Serie}

La zona intermedia está aislada: su carga neta es cero, forzando la misma
$Q$ en ambos. Las diferencias de potencial se suman:
\[
\Delta V_{AB} = \Delta V_{AK} + \Delta V_{KB} = -\frac{Q}{C_1} - \frac{Q}{C_2}.
\]

\begin{keybox}[Serie]
  \[
  \boxed{\frac{1}{C_{\text{eq}}} = \frac{1}{C_1} + \frac{1}{C_2}}
  \]
\end{keybox}

\begin{notebox}[Consecuencia]
  En serie $C_{\text{eq}} \leq C_1$ (y $\leq C_2$): la capacidad
  equivalente es \textbf{menor} que cualquiera de las individuales.
\end{notebox}

\subsection{Ejemplo combinado}

$C_1$ y $C_2$ en paralelo, en serie con $C_3$. Por etapas:
$C_{12} = C_1 + C_2$, luego
\[
\frac{1}{C_{\text{eq}}} = \frac{1}{C_1 + C_2} + \frac{1}{C_3}.
\]

\begin{warnbox}[No siempre se puede]
  Si hay baterías, resistencias u otros elementos entre medio, este método
  de reducción no aplica.
\end{warnbox}

\subsection{Uso práctico}

Para ajustar una capacidad no disponible: \textbf{paralelo} para
aumentar, \textbf{serie} para disminuir.

% ===================================================================
\section{Energía almacenada en un condensador}
% ===================================================================

Con carga $Q'$ ($\Delta V' = Q'/C$), agregar $dQ'$ aumenta la energía:
\[
dU = \Delta V'\, dQ' = \frac{Q'}{C}\,dQ'.
\]

Con $U=0$ descargado, integrando:

\begin{keybox}[Energía de un condensador]
  \[
  \boxed{U = \frac{Q^2}{2C} = \frac{1}{2}C\,\Delta V^2 = \frac{1}{2}Q\,\Delta V}
  \]
\end{keybox}

% ===================================================================
\section{Energía almacenada en el campo eléctrico}
% ===================================================================

Para placas paralelas ($C = \varepsilon_0 A/d$, $\Delta V = E d$):
\[
U = \frac{1}{2}C\,\Delta V^2
  = \frac{1}{2}\frac{\varepsilon_0 A}{d}(E d)^2
  = \frac{1}{2}\varepsilon_0 E^2\,(A d),
\]
donde $A d$ es el volumen interior. La energía está distribuida en el
campo, con densidad volumétrica:

\begin{keybox}[Densidad de energía del campo eléctrico]
  \[
  \boxed{u = \frac{1}{2}\varepsilon_0 E^2}, \qquad U = \int u\, dV
  \]
\end{keybox}

\begin{notebox}[Validez]
  Deducida solo para placas paralelas, pero vale siempre en el vacío
  (demostración en Electromagnetismo). Para la energía de un sistema:
  integrar $u$ sobre toda la región donde $\mathbf{E}\neq 0$.
\end{notebox}

% ===================================================================
\section{Anuncio: aislante en un campo (dieléctricos)}
% ===================================================================

¿Qué pasa si se rellena el condensador con un \textbf{aislante}
(dieléctrico)? No puede ser conductor (descargaría el capacitor).

\textbf{Polarización:} un aislante de moléculas polares responde al
campo:
\begin{itemize}[nosep]
  \item $Q=0$: moléculas orientadas al azar.
  \item $Q\neq 0$: se orientan \textbf{levemente} en promedio (fuerzas
        internas del material se oponen a una alineación total).
\end{itemize}

\textbf{Efecto neto:} \textbf{cargas de polarización} en las superficies
del aislante (negativas junto a la placa positiva, positivas junto a la
negativa). El interior sigue neutro; las cargas no se desplazan (no hay
conducción).

\begin{warnbox}[Dificultad]
  Normalmente no se sabe cuánto se polariza un material (problema de
  \textbf{mecánica estadística}). La próxima clase: cómo describir el
  dieléctrico sin calcular la polarización exacta.
\end{warnbox}

\end{document}
